% Reconstructed after environment loss from accepted proof and retained conversation, 16 November 2026.
\statement{C.41}{Proved}{An explicit first relation-module bound for the canonical word}
Let $F=\widehat F(x_1,\ldots,x_n)$, let $c(x_i)=x_i^{-1}$,
and let $q:F\to Q$ be a finite quotient with $c$-stable kernel $U$.
Put $\Phi_2(U)=\overline{U^2[U,U]}$. Then
\[
 C_{F/\Phi_2(U)}(c)\longrightarrow F/U
                 \quad\hbox{has trivial image}.                 \tag{C.41.1}
\]
For the rank-two canonical equation $u=f\alpha(u)$,
where $\alpha(x)=x^{-1}$ and $\alpha(y)=f^{-1}y^{-1}f$,
this gives the explicit implication
\[
                  f\in\Phi_2(U)\quad\Longrightarrow\quad u\in U.
                                                                    \tag{C.41.2}
\]
The latter assertion concerns the actual full canonical word and a
$c$-stable $U$. It does not say that the fixed subgroup inside the
relation module itself is trivial.

\begin{proof}
Write $a_i=q(x_i)$ and let $\kappa$ be the induced involution of $Q$.
Consider the finite $Q$-cover of the rose with $n$ oriented loops.
Its vertices are $g\in Q$, and its positively oriented $i$th edge at
$g$ joins $g$ to $ga_i$. Work with its mod-two chain modules
$C_1,C_0$. There is an actual finite quotient model
\[
 F/\Phi_2(U)\cong
 \{(v,g)\in C_1\rtimes Q:\partial v=e_g+e_1\},                  \tag{C.41.3}
\]
where $Q$ acts by left translation and $x_i$ maps to its edge at $1$
together with $a_i$.

To check this model without replacing the full topology, let $F_0$
be the discrete free group and $U_0=U\cap F_0$. The edge chain of
a word is its actual lifted path. For a fixed endpoint, the difference
of two such chains is a cycle. Every mod-two cycle is the homology
class of a based closed path in this connected graph. Thus the image
is precisely the group on the right of (C.41.3), and the kernel on
$U_0$ is $U_0^2[U_0,U_0]$. Finite-index topology identifies
$\widehat{U_0}$ with $U$; its finite mod-two homology identifies the
closure of this kernel with $\Phi_2(U)$. This proves (C.41.3) as a
full profinite quotient.

The involution sends vertex $g$ to $\kappa(g)$ and the positive
$i$th edge at $g$ to the reversed $i$th edge at
$\kappa(g)a_i^{-1}$. In mod-two chains its reversal sign disappears.
Let $v$ be an invariant edge chain. At a $\kappa$-fixed vertex, a
pair of exchanged edges has boundary coefficient zero: the two
coefficients are equal. A fixed edge is reversed and therefore has
its endpoints exchanged. It cannot meet a fixed vertex unless it
is a loop, whose boundary is zero. Consequently
\[
 (\partial v)(g)=0\quad\hbox{at every $\kappa$-fixed vertex $g$}.
                                                                    \tag{C.41.4}
\]
If $(v,g)$ is fixed and $g\ne1$, then $\kappa(g)=g$ and its
boundary $e_g+e_1$ has coefficient one at the fixed vertex $1$.
This contradicts (C.41.4), proving (C.41.1).

If $f\in\Phi_2(U)$, the full automorphism $\alpha$ induces $c$
on (C.41.3). This also proves the required stability of the kernel:
first $q\alpha=\kappa q$ makes $U$ stable, and then $\Phi_2(U)$
is characteristic in $U$. The actual canonical equation gives a
$c$-fixed image of $u$, and (C.41.1) makes its $Q$-image one.
\end{proof}

\paragraph{Exact finite certificate and scope.}
The reconstructed, rerun certificate
\path{work_CE/code/check_C41_relation_module.py}
uses the critical $(3,5,2)$ icosahedral head. Its graph has $60$
vertices, $120$ edges and cycle rank $61$. The invariant-chain
space has dimension $64$, its boundary image has rank $28$, and
its invariant cycle space has dimension $36$. All three
nonidentity fixed head elements are excluded by (C.41.4).
These are checks of the explicit finite model; the proof above
works in every rank and every finite head.
The earlier false-positive $A_5$ table specified only $f_Q=1$.
It can still have a nonzero parameter in this relation module,
so that table is excluded here only under the additional hypothesis
that $f$ is trivial in the enlarged quotient.

Iteration also recovers the relative pro-two fixed-group assertion
already proved in A.29. For $c$-stable open normal $V\subset U$
with $U/V$ a finite two-group, one has $O^2(V)=O^2(U)$: the quotient
$U/O^2(V)$ is pro-two, and the reverse inclusion follows by
restriction of pro-two quotients. The next quotient
$F/\Phi_2(V)$ therefore lies in the same relative tower and kills
every possible nonidentity fixed finite head. The new content of
C.41 is its explicit one-step bound, not a new attribution of A.29.
